%% These are LaTeX dependencies that are needed to build the paper but are not a part of the conference template

% \usepackage{usenix}
\usepackage{usenix2021}
\usepackage{filecontents}
\usepackage{enumitem}
\usepackage{subcaption}
\usepackage{graphicx}
\usepackage{latexsym}
\usepackage{amsmath}
\usepackage{listings}
\usepackage{booktabs}
\usepackage{amssymb}
\usepackage{pstricks}
\usepackage{comment}
\usepackage{xspace}
\usepackage{multirow}
\usepackage{makecell}
\usepackage{tabularx}
\usepackage{tikz}
%\usepackage{pifont}
\usepackage{wrapfig}
\usepackage[hyphens]{xurl}
% \usepackage{breakcites}

%\usepackage{cite}
%\usepackage[numbers,super]{natbib}
\usepackage{xcolor}
\usepackage{amsfonts}

% \usepackage{adjustbox}
\usepackage{tikz}
\usepackage[available,functional,reproduced]{usenixbadges}
\usetikzlibrary{arrows,tikzmark}

% Set title/author height
% \usepackage{titling}
% \setlength{\droptitle}{-2.5em}
% \postauthor{\end{tabular}\vspace{-0.5in}\end{center}}
% \date{}

% Paragraph and caption spacing
% \setlength{\abovecaptionskip}{5pt}
% \setlength{\belowcaptionskip}{10pt}
% \setlength{\floatsep}{3pt}
% \setlength{\textfloatsep}{3pt}
% \setlength{\dbltextfloatsep}{3pt}
% \setlength{\intextsep}{3pt}

% \usepackage{sectsty}
% \subsectionfont{\fontsize{12}{16}\selectfont}
% \subsubsectionfont{\fontsize{11}{14}\selectfont}
% \paragraphfont{\fontsize{10.5}{12.5}\selectfont}

\renewcommand*\ttdefault{txtt}

% Set title/author height
\usepackage{titling}
\setlength{\droptitle}{-3.05em}
\postauthor{\end{tabular}\vspace{-0.5in}\end{center}}
\date{}

% Use these to save space
\setlength{\abovecaptionskip}{5pt}
\setlength{\belowcaptionskip}{10pt}
\setlength{\floatsep}{2.5pt}
\setlength{\textfloatsep}{3pt}
\setlength{\dbltextfloatsep}{3pt}
\setlength{\intextsep}{3pt}

% Section header spacing
\usepackage{titlesec}
\titlespacing\subsection{0pt}{12.5pt}{4.5pt}
\titlespacing\subsubsection{0pt}{12pt}{3pt}
\titlespacing\paragraph{0pt}{12pt}{3pt}
\titleformat{\subsection}
  {\normalfont\fontsize{12}{15}\bfseries}{\thesection}{1em}{}
\titleformat{\subsubsection}
  {\normalfont\fontsize{11}{15}\bfseries}{\thesubsection}{1em}{}

\usepackage{minted}
\newcommand{\codehl}[3]{{\setlength{\fboxsep}{1pt}\textbf{\colorbox{hlgray}{#2\tikzmark{#1}#3}}}}
\newcommand{\codenote}[6]{\draw[-,overlay,thin] (pic cs:#1) +(0,#2) |- +(#3,#4) node[rectangle,draw,#5,thin,fill=bkgray,inner sep=2pt] {\scriptsize\makecell[l]{#6}};}
\newcommand{\codetype}[1]{\textbf{\color{mintedtype}#1}}
\newcommand{\codefunc}[1]{\color{mintedfunc}#1}
\newcommand{\coderet}[1]{\textbf{\color{mintedret}#1}}

% \renewcommand*{\ttdefault}{cmtt}

\newenvironment{packed_itemize}{
    \begin{list}{\labelitemi}{\leftmargin=1.0em}
     \setlength{\itemsep}{2.5pt}
     \setlength{\parskip}{0pt}
     \setlength{\parsep}{0pt}
     \setlength{\headsep}{0pt}
     \setlength{\topskip}{0pt}
     \setlength{\topmargin}{0pt}
     \setlength{\topsep}{0pt}
     \setlength{\partopsep}{0pt}
    }{\end{list}}

\newcommand*\circled[1]{%
	\begin{tikzpicture}[baseline=(C.base)]
		\node[draw,circle,fill=black,inner sep=0.1pt](C) {\textcolor{white}{#1}};
	\end{tikzpicture}}

\newcommand{\aspy}{\makecell{\checkmark}}
\newcommand{\aspn}{\makecell{-}}

\definecolor{hlgray}{HTML}{DDDDDD}
\definecolor{bkgray}{HTML}{EEEEEE}
\definecolor{mintedtype}{HTML}{B00040}
\definecolor{mintedhl}{HTML}{C00000}
\definecolor{mintedfunc}{HTML}{0000FF}
\definecolor{mintedret}{HTML}{999999}

\definecolor{comment}{HTML}{008000}

\definecolor{popipa}{HTML}{FF3377}
\definecolor{roselia}{HTML}{3344AA}
\definecolor{ras}{HTML}{39C9C5}
\definecolor{monica}{HTML}{33AAFF}
\definecolor{mygo}{HTML}{3388BB}
\definecolor{mujica}{HTML}{881144}

\definecolor{diffadd}{HTML}{80E597}
\definecolor{diffdel}{HTML}{FF9F99}
\definecolor{diffoth}{HTML}{AAAAAA}
\definecolor{diffamv}{HTML}{DAFBE1}
\definecolor{diffdmv}{HTML}{FFEBE9}

\newcommand{\TODO}[1]{{\color{red} #1}}
\newcommand{\REVIEW}[1]{{\color{blue} #1}}


% Editorial
\iftrue % change to \iffalse to preview changes

	\usepackage[normalem]{ulem}

	\newcommand{\JZ}[1]{{\color{orange} #1}} % new text
	\newcommand{\JZx}[1]{{\color{purple} \sout{#1}}} % del text
	\newcommand{\JZc}[1]{{\color{comment} /* \textit{#1} */}} % comment
	\newcommand{\JZd}[1]{} % exile to the void
	\newcommand{\JZe}[1]{{\color{mujica} \uwave{#1}}} % highlight
	\newcommand{\JZhl}{\colorbox{mujica}{\color{white} NOTE:}}
	\newcommand{\JZtodo}{\colorbox{mygo}{\color{white} TODO:}}

	\newcommand{\jz}[1]{{\JZ{#1}}} % new text
	\newcommand{\jza}[2]{{\JZ{#1} \JZc{#2}}} % add and comment text
	\newcommand{\jzc}[2]{{\JZe{#1} \JZc{#2}}} % comment existing text
	\newcommand{\jzd}[2]{{\JZx{#1} \JZc{#2}}} % delete and comment text
	\newcommand{\jzx}[2]{{\JZx{#1} \JZ{#2}}} % delete and replace text
	\newcommand{\jzxc}[3]{{\JZx{#1} \JZ{#2} \JZc{#3}}} % replace text with comment
	\newcommand{\jztodo}[1]{{\JZc{\JZtodo #1}}}
	\newcommand{\jzmark}[1]{{\JZc{\JZhl #1}}}


	% \newcommand{\schai}[1]{{\color{blue} #1}} % new text
	\newcommand{\schai}[1]{#1}
	\newcommand{\schaix}[1]{{\color{red} \sout{#1}}} % del text
	\newcommand{\schaic}[1]{{\color{comment} /* \bf schai \textit{#1} */}} % comment
	
	\newcommand{\alan}[1]{{\color{comment} /* \bf alan \textit{#1} */}} % comment

	\newcommand{\jyk}[1]{{\footnotesize{\color{cyan} {\bf JYK: #1}}}}
	\newcommand{\jykx}[1]{{\color{cyan} \sout{#1}}}

	\newcommand{\weiwei}[1]{{\color{olive} {\bf WEIWEI: #1}}}
	\newcommand{\weiweix}[1]{{\color{red} \sout{#1}}}

	% so we can mix \cite within colored text
	\let\oldcite\cite
	\renewcommand{\cite}[1]{\mbox{\oldcite{#1}}}

\else
	\newcommand{\tianyin}[1]{#1}

	\newcommand{\JZ}[1]{#1} % new text
	\newcommand{\JZx}[1]{} % del text
	\newcommand{\JZc}[1]{} % comment
	\newcommand{\JZd}[1]{} % exile to the void
	\newcommand{\JZe}[1]{} % highlight
	\newcommand{\JZhl}{}
	\newcommand{\JZtodo}{}

	\newcommand{\jz}[1]{#1}
	\newcommand{\jzd}[2]{#2}
	\newcommand{\jzc}[2]{#1}
	\newcommand{\jzx}[2]{#2}
	\newcommand{\jzxc}[3]{#2}
	\newcommand{\jztodo}{}

	\newcommand{\schai}[1]{#1} % new text
	\newcommand{\schaix}[1]{} % del text
	\newcommand{\schaic}[1]{} % comment

	\newcommand{\jyk}[1]{}
	\newcommand{\jykx}[1]{}

	\newcommand{\weiwei}[1]{}
	\newcommand{\weiweix}[1]{}
\fi

% Layout
\usepackage{pgffor}
\newcommand{\para}[1]{\vspace{.10in}\noindent{\bf #1}}
\newcommand{\ident}[1]{\foreach \index in {1, ..., #1} { }}
\newcommand{\diffline}[2]{\setlength{\fboxsep}{3pt}\colorbox{#1}{\raisebox{0pt}[0.5em][0pt]{\makebox[\textwidth][l]{\!\!#2}}}}
\newcommand{\diffadd}[1]{\diffline{diffadd}{#1}}
\newcommand{\diffdel}[1]{\diffline{diffdel}{#1}}
\newcommand{\diffoth}[1]{\diffline{diffoth}{#1}}
\newcommand{\diffamv}[1]{\diffline{diffamv}{#1}}
\newcommand{\diffdmv}[1]{\diffline{diffdmv}{#1}}

% Paper Specific
\newcommand{\branding}{EMT\xspace} % Stands for Extensible Memory Translation
\newcommand{\interface}{\branding interface\xspace}

\newcommand{\eradix}{\branding-Radix\xspace}
\newcommand{\eecpt}{\branding-ECPT\xspace}

\newcommand{\iigrp}{Functional Group\xspace}
\newcommand{\iigrps}{function groups\xspace}
\newcommand{\igrp}{\branding functional group\xspace}
\newcommand{\igrps}{\branding functional groups\xspace}

\newcommand{\ptsfg}{{\it PTS-fg}\xspace}
\newcommand{\ptefg}{{\it PTE-fg}\xspace}
\newcommand{\hwofg}{{\it HWO-fg}\xspace}

\newcommand{\iilyr}{Abstraction Layer\xspace}
\newcommand{\ilyr}{\branding abstraction layer\xspace}
\newcommand{\ilyrs}{\branding abstraction layers\xspace}

\newcommand{\mmudriver}{MMU driver\xspace}
\newcommand{\mmudrivers}{MMU drivers\xspace}
\newcommand{\essgrps}{Primitives\xspace}
\newcommand{\optgrps}{Optimizations\xspace}

\newcommand{\pte}{PSE\xspace}
\newcommand{\ptes}{PSEs\xspace}
\newcommand{\ptefull}{paging structure entry\xspace}
\newcommand{\ptesfull}{paging structure entries\xspace}

\usepackage{array}
\newcolumntype{F}[1]{>{\arraybackslash}p{#1}}
\newcommand{\itblhead}[1]{\hline \multicolumn{10}{|c|}{\textbf{#1}} \\ \hline \hline}
\newcommand{\itblnext}[1]{\hline \hline \multicolumn{10}{|c|}{\textbf{#1}} \\ \hline \hline}
\newcommand{\itblitem}[2]{\multicolumn{10}{|l|}{\texttt{\makecell[l]{#1}}} \\ & \multicolumn{9}{F{0.91\linewidth}|}{#2} \\ \hline}
\newcommand{\itblnote}[1]{\multicolumn{10}{|l|}{#1} \\ \hline}

% Compact list
\usepackage{paralist}

% Subfigure
\usepackage{subcaption}